\section{From spherical angles to rational orientation charts}
\label{cp:sec-phase-charts}

This section replaces the trigonometric Regge transformation by an
algebraic family.  There are three steps.  First, exponentiating the
dihedral angles turns the linear Regge transformation into a monomial
involution of an algebraic torus.  Second, adjoining a square root of the
Gram determinant retains the orientation of the associated quadric.  Third,
a change of variable rationalizes this double cover.  The only angle phase
which is lost in the rationalization is recovered by one explicit quadratic
finite \'etale cover.

No motivic input is used in this section.

\subsection{The spherical angle Gram matrix}
\label{cp:subsec-angle-gram}

Let \(T\subset S^3\) be a nondegenerate spherical tetrahedron contained in
an open hemisphere.  Write its four facets as \(F_0,\ldots,F_3\), and let
\(n_i\in S^3\subset\mathbf R^4\) be the outward unit normal to \(F_i\).
If \(F_i\) and \(F_j\) meet at interior dihedral angle
\(\theta_{ij}\in(0,\pi)\), then their outward normals meet at the
supplementary angle.  Hence
\begin{equation}
 n_i\mathbin{\cdot}n_j=-\cos\theta_{ij}.
 \label{cp:eq-normal-angle}
\end{equation}
The geometric edge at which these facets meet is indexed by the
complementary pair of vertices.  We use the facet-pair labels \(ij\), since
they are the natural labels for the Gram matrix.  Complementing all six
labels preserves opposite pairs and merely changes the edge-label
convention.

\begin{definition}[Angle Gram matrix]
\label{cp:def-angle-gram}
The angle Gram matrix of \(T\) is
\begin{equation}
 G(T)=(g_{ij})_{0\leq i,j\leq3},\qquad
 g_{ii}=1,\qquad g_{ij}=-\cos\theta_{ij}\quad(i\ne j).
 \label{cp:eq-angle-gram}
\end{equation}
\end{definition}

If \(N\) is the matrix with columns \(n_0,\ldots,n_3\), then
\(G=N^tN\).  In particular, a nondegenerate spherical tetrahedron has
\begin{equation}
 G>0,\qquad \det G=(\det N)^2>0.
 \label{cp:eq-spherical-gram-positive}
\end{equation}
The positivity is used only to locate the physical real locus.  All
algebraic constructions below allow the entries of \(G\) to be complex.

Fix the opposite Gram positions \(01\) and \(23\).  In the cyclic order
\begin{equation}
 (01,02,12,23,13,03),
 \label{cp:eq-edge-order}
\end{equation}
write the two fixed angles as \(\theta_{01},\theta_{23}\) and the four
moving angles as
\begin{equation}
 \boldsymbol\theta
   =(\theta_{02},\theta_{12},\theta_{13},\theta_{03})^t.
 \label{cp:eq-moving-angle-vector}
\end{equation}
Put \(\mathbf 1=(1,1,1,1)^t\), let
\(J_4=\mathbf1\mathbf1^t\), and define
\begin{equation}
 R=\frac12J_4-I_4.
 \label{cp:eq-regge-matrix}
\end{equation}
If
\begin{equation}
 S=\frac12(\theta_{02}+\theta_{12}+\theta_{13}+\theta_{03}),
 \label{cp:eq-regge-semisum}
\end{equation}
then
\begin{equation}
 R\boldsymbol\theta
 =
 (S-\theta_{02},S-\theta_{12},S-\theta_{13},S-\theta_{03})^t.
 \label{cp:eq-linear-regge-action}
\end{equation}

\begin{lemma}[Elementary properties of the Regge matrix]
\label{cp:lem-regge-matrix}
The matrix \(R\) is a symmetric orthogonal involution:
\[
 R^t=R,\qquad R^2=I_4.
\]
It is the identity on \(\mathbf Q\mathbf1\) and minus the identity on
\(\mathbf1^\perp\).
\end{lemma}

\begin{proof}
Since \(J_4^2=4J_4\),
\[
 R^2=\frac14J_4^2-J_4+I_4=I_4.
\]
Symmetry is immediate.  Also \(R\mathbf1=\mathbf1\), whereas
\(J_4v=0\) and hence \(Rv=-v\) for \(v\in\mathbf1^\perp\).
\end{proof}

\begin{theorem}[Geometric Regge input]
\label{cp:thm-regge-input}
An elementary Regge move fixes \(\theta_{01},\theta_{23}\) and replaces
\(\boldsymbol\theta\) by \(R\boldsymbol\theta\).  The same matrix acts on
the corresponding four edge lengths, and the two tetrahedra have equal
volume.
\end{theorem}

\begin{proof}
This is the constant-curvature Regge theorem of Akopyan and Izmestiev
\cite[Theorem~1]{AkopyanIzmestiev2019}.  Only the angle statement is needed
in the present section.
\end{proof}

\subsection{Phase coordinates and the meaning of ``monomial''}
\label{cp:subsec-phase-coordinates}

For a physical tetrahedron define its six angle phases by
\begin{equation}
 \begin{aligned}
 p&=e^{i\theta_{01}},&
 a&=e^{i\theta_{02}},&
 b&=e^{i\theta_{12}},\\
 q&=e^{i\theta_{23}},&
 c&=e^{i\theta_{13}},&
 d&=e^{i\theta_{03}},
 \end{aligned}
 \qquad \tau=e^{iS}.
 \label{cp:eq-angle-phases}
\end{equation}
Then
\begin{equation}
 \tau^2=abcd.
 \label{cp:eq-phase-relation}
\end{equation}
Exponentiating \eqref{cp:eq-linear-regge-action} gives
\begin{equation}
 (p,a,b,q,c,d,\tau)\longmapsto
 \left(p,\frac\tau a,\frac\tau b,q,
              \frac\tau c,\frac\tau d,\tau\right).
 \label{cp:eq-monomial-regge-action}
\end{equation}
This is the exact meaning of the phrase ``the Regge transformation acts
monomially.''  A Laurent monomial in variables \(z_1,\ldots,z_m\) is a
product \(z_1^{n_1}\cdots z_m^{n_m}\) with integral exponents.  Every new
coordinate in \eqref{cp:eq-monomial-regge-action} is such a Laurent
monomial.  The additive formula \(S-\theta\) has become the multiplicative
formula \(\tau/a\).  Thus no choice of logarithm occurs in the algebraic
formula.

To make that observation scheme-theoretic, set
\begin{equation}
 \mathcal X=
 \operatorname{Spec}
 \frac{\mathbf Q[p^{\pm1},a^{\pm1},b^{\pm1},q^{\pm1},
                         c^{\pm1},d^{\pm1},\tau^{\pm1}]}
      {(\tau^2-abcd)}.
 \label{cp:eq-phase-torus}
\end{equation}
Here \(\Spec R\) means the affine algebraic space whose regular functions
form the ring \(R\), and \(\mathbf G_m=\Spec\Q[u,u^{-1}]\) is the affine
line with the point \(u=0\) removed.  Thus an \(F\)-point of
\(\mathcal X\), for a field \(F\supset\Q\), is simply a tuple of nonzero
elements of \(F\) satisfying \(\tau^2=abcd\).  This notation packages
polynomial identities and changes of coordinates; no extra geometric
hypothesis is hidden in it.
Eliminating \(d=\tau^2/(abc)\) identifies \(\mathcal X\) with the
six-dimensional algebraic torus \(\mathbf G_m^6\).  Physical phase data form
a real-analytic subset of \(\mathcal X(\mathbf C)\), characterized in part
by \(|p|=|a|=\cdots=|q|=1\).  We enlarge to all complex points because
\(\mathcal X\), unlike the physical locus, is an algebraic parameter space
defined over \(\mathbf Q\).

\begin{lemma}[Algebraic Regge involution]
\label{cp:lem-algebraic-regge-involution}
Formula \eqref{cp:eq-monomial-regge-action} defines an involutive
automorphism \(\mathcal R:\mathcal X\to\mathcal X\).
\end{lemma}

\begin{proof}
The product of the four transformed moving phases is
\[
 \frac{\tau^4}{abcd}=\tau^2,
\]
so the defining relation of \(\mathcal X\) is preserved.  Applying the map
twice sends, for example, \(a\) to \(\tau/(\tau/a)=a\), and similarly for
the other coordinates.  The fixed coordinates remain fixed.
\end{proof}

The phase \(u=e^{i\theta}\) recovers the corresponding Gram entry without
choosing \(\theta\):
\begin{equation}
 -\cos\theta=-\frac{u+u^{-1}}2.
 \label{cp:eq-phase-to-cosine}
\end{equation}
Accordingly put
\begin{equation}
 \begin{gathered}
 x=-\frac{p+p^{-1}}2,\qquad
 A=-\frac{a+a^{-1}}2,\qquad
 B=-\frac{b+b^{-1}}2,\\
 y=-\frac{q+q^{-1}}2,\qquad
 C=-\frac{c+c^{-1}}2,\qquad
 D=-\frac{d+d^{-1}}2.
 \end{gathered}
 \label{cp:eq-phase-gram-entries}
\end{equation}
In the order fixed above, the universal angle Gram matrix is
\begin{equation}
 G=
 \begin{pmatrix}
  1&x&A&D\\
  x&1&B&C\\
  A&B&1&y\\
  D&C&y&1
 \end{pmatrix}.
 \label{cp:eq-gram}
\end{equation}
At a physical point this specializes to \eqref{cp:eq-angle-gram}.

\subsection{The determinant orientation cover}
\label{cp:subsec-orientation-cover}

On the open subset of \(\mathcal X\) where \(\det G\ne0\), adjoin a
variable \(w\) satisfying
\begin{equation}
 w^2=\det G.
 \label{cp:eq-orientation-root}
\end{equation}
Because the characteristic is zero and \(w\ne0\), this is a finite
\'{e}tale double cover.  Its deck involution is \(w\mapsto-w\).
A finite \'{e}tale map is the algebraic analogue of a finite covering
without branch points.  Here this is elementary: the defining polynomial is
\(w^2-\det G\), and its derivative \(2w\) is invertible on the chosen open,
so its two roots never collide.

The terminology ``orientation cover'' refers to the projective quadric
which will be attached to \(G\) later.  A smooth quadric surface over an
algebraically closed field has two families of maximal isotropic lines.  For
a four-dimensional quadratic form, choosing a square root of its
determinant distinguishes these two families; changing the sign of the root
interchanges them.  Thus \(w\), rather than the matrix \(G\) alone, retains
the orientation needed by the quadric-scissors construction.  At a physical
spherical point we use the positive real root \(w=\sqrt{\det G}\).

The determinant is Regge invariant.  We now prove this in a form which also
rationalizes the orientation cover.

Keep \(a,b,c,\tau,q\) as independent invertible variables and eliminate
\begin{equation}
 d=\frac{\tau^2}{abc}.
 \label{cp:eq-eliminate-d}
\end{equation}
We temporarily forget the phase \(p\), retaining only its cosine coordinate
\(x\).  Define
\begin{equation}
 s=\frac{q-q^{-1}}2,
 \qquad y=-\frac{q+q^{-1}}2.
 \label{cp:eq-q-and-s}
\end{equation}
Then
\begin{equation}
 1-y^2=-s^2.
 \label{cp:eq-y-s-relation}
\end{equation}
The asymmetry between \(p\) and \(q\) is deliberate: we descend \(p\) to
\(x\), while the retained phase \(q\) supplies the square root \(s\) in
\eqref{cp:eq-y-s-relation}.

Define
\begin{align}
 \mathcal B&=2\{AB+DC-y(AC+DB)\},
 \label{cp:eq-determinant-B}\\
 \mathcal C&=1-A^2-B^2-C^2-D^2-y^2
       +2y(AD+BC)+(AC-BD)^2.
 \label{cp:eq-determinant-C}
\end{align}

\begin{lemma}[Determinant quadratic and Regge invariance]
\label{cp:lem-determinant-quadratic}
For the matrix \eqref{cp:eq-gram},
\begin{equation}
 \det G=s^2x^2+\mathcal Bx+\mathcal C.
 \label{cp:eq-determinant-quadratic}
\end{equation}
Under the monomial Regge substitution
\begin{equation}
 (a,b,c,d)\longmapsto
 \left(\frac\tau a,\frac\tau b,
       \frac\tau c,\frac\tau d\right),
 \label{cp:eq-moving-phase-regge}
\end{equation}
with \(p,q,\tau\) fixed, one has
\begin{equation}
 \mathcal B^R=\mathcal B,\qquad
 \mathcal C^R=\mathcal C.
 \label{cp:eq-determinant-coefficient-invariance}
\end{equation}
Consequently \(\det G^R=\det G\), and the Regge involution has the lift
\begin{equation}
 (G,w)\longmapsto(G^R,w)
 \label{cp:eq-oriented-regge-lift}
\end{equation}
which does not change the chosen determinant root.
\end{lemma}

\begin{proof}
Expand the determinant of \eqref{cp:eq-gram} as a polynomial
in \(x\).  Its quadratic coefficient is
\(-(1-y^2)=s^2\), its linear coefficient is
\(2\{AB+DC-y(AC+DB)\}\), and its constant coefficient is the expression in
\eqref{cp:eq-determinant-C}.  This proves
\eqref{cp:eq-determinant-quadratic}.

For the invariance assertion, put
\[
 h(u)=\frac{u+u^{-1}}2
\]
and introduce
\begin{align*}
 I_1&=AB+CD,& I_2&=AC+BD,& I_3&=AD+BC,\\
 K&=A^2+B^2+C^2+D^2+4ABCD.
\end{align*}
The elementary product-to-sum identity
\[
 2h(u)h(v)=h(uv)+h(u/v)
\]
is an identity of Laurent polynomials.  Under
\(a'=\tau/a,\ldots,d'=\tau/d\), the relation \(abcd=\tau^2\) gives
\[
 a'b'=cd,\quad c'd'=ab,\quad
 a'/b'=b/a,\quad c'/d'=d/c.
\]
Since \(h(u^{-1})=h(u)\), the product-to-sum identity shows that
\(I_1^R=I_1\).  Applying the same argument to the partitions
\(ac\mid bd\) and \(ad\mid bc\) gives \(I_2^R=I_2\) and \(I_3^R=I_3\).

It remains to treat \(K\).  Set
\[
 P=h(ab),\quad Q=h(cd),\quad r=h(a/b),\quad s_0=h(c/d).
\]
Direct applications of the same product-to-sum identity give
\[
\begin{aligned}
 h(a)^2+h(b)^2&=1+Pr,\\
 h(c)^2+h(d)^2&=1+Qs_0,\\
 4h(a)h(b)h(c)h(d)&=(P+r)(Q+s_0).
\end{aligned}
\]
Because \(A=-h(a),\ldots,D=-h(d)\), it follows that
\[
 K=2+PQ+rs_0+(P+Q)(r+s_0).
\]
The Regge substitution swaps \(P,Q\) and fixes \(r,s_0\), so \(K^R=K\).
Finally,
\[
 \frac{\mathcal B}{2}=I_1-yI_2,\qquad
 \mathcal C=1-y^2+2yI_3+I_2^2-K,
\]
where the second equality uses
\((AC-BD)^2=(AC+BD)^2-4ABCD\).  Thus
\(\mathcal B^R=\mathcal B\) and \(\mathcal C^R=\mathcal C\) inside the
Laurent coordinate ring of \(\mathcal X\), not merely on the physical
locus.  Since \(x,y,s\) are fixed, the determinant quadratics coincide.
Keeping the same root \(w\) therefore gives the displayed lift.
\end{proof}

\subsection{Rational charts on the orientation cover}
\label{cp:subsec-rational-charts}

Let
\[
 S_0=\mathbf Q[a^{\pm1},b^{\pm1},c^{\pm1},
                       \tau^{\pm1},q^{\pm1}],
\]
and regard \(d,A,B,C,D,y,s,\mathcal B,\mathcal C\) as the elements of
\(S_0\) defined above.  The descended orientation space is the affine
hypersurface
\begin{equation}
 \mathcal Y=
 \operatorname{Spec}
 \frac{S_0[x,w]}
      {(w^2-s^2x^2-\mathcal Bx-\mathcal C)}.
 \label{cp:eq-descended-orientation-space}
\end{equation}
The orientation cover of the phase torus maps to \(\mathcal Y\) by
\(x=-(p+p^{-1})/2\).  The point of passing to \(\mathcal Y\) is that it is
rational, although the original equation visibly contains a square root.

\begin{proposition}[Two rational orientation charts]
\label{cp:lem-rational-charts}
Put
\begin{equation}
 r=w-sx,\qquad r_+=w+sx.
 \label{cp:eq-rational-chart-coordinates}
\end{equation}
On the principal open set where \(2sr-\mathcal B\ne0\), one has
\begin{equation}
 x=\frac{\mathcal C-r^2}{2sr-\mathcal B},
 \qquad w=r+sx.
 \label{cp:eq-minus-rational-chart}
\end{equation}
On the principal open set where \(-2sr_+-\mathcal B\ne0\), one has
\begin{equation}
 x=\frac{\mathcal C-r_+^2}{-2sr_+-\mathcal B},
 \qquad w=r_+-sx.
 \label{cp:eq-plus-rational-chart}
\end{equation}
Each formula identifies the indicated part of \(\mathcal Y\) with an open
subset of \(\mathbf G_m^5\times\mathbf A^1\).  In particular, its function
field is purely transcendental:
\begin{equation}
 \mathbf Q(a,b,c,\tau,q,r)
 \quad\text{or}\quad
 \mathbf Q(a,b,c,\tau,q,r_+).
 \label{cp:eq-chart-function-fields}
\end{equation}
On the locus \(sw\ne0\), the two charts cover \(\mathcal Y\).
\end{proposition}

\begin{proof}
Substitute \(w=r+sx\) into
\eqref{cp:eq-determinant-quadratic}:
\[
 r^2+2rsx+s^2x^2=s^2x^2+\mathcal Bx+\mathcal C.
\]
After cancellation of the quadratic terms,
\[
 (2sr-\mathcal B)x=\mathcal C-r^2,
\]
which is \eqref{cp:eq-minus-rational-chart}.  Conversely, the two formulas
in \eqref{cp:eq-minus-rational-chart} satisfy both
\(r=w-sx\) and the defining equation of \(\mathcal Y\); hence they give
mutually inverse coordinate maps on the stated principal opens.  The proof
of \eqref{cp:eq-plus-rational-chart} is identical after substituting
\(w=r_+-sx\).

It remains to check coverage.  Written in the original coordinates, the
two denominators are
\[
 2s(w-sx)-\mathcal B,\qquad
 -2s(w+sx)-\mathcal B.
\]
If both vanished, subtracting the second equation from the first would give
\(4sw=0\).  This is impossible on the locus \(sw\ne0\).
\end{proof}

The change of variable has a simple geometric-algebraic interpretation.
The equation \(w^2=s^2x^2+\mathcal Bx+\mathcal C\) is a conic in the
\((x,w)\)-plane over the field
\(\mathbf Q(a,b,c,\tau,q)\).  The coordinates \(r=w-sx\) and
\(r_+=w+sx\) subtract one or the other asymptotic linear part.  This cancels
the \(x^2\)-term and leaves an equation linear in \(x\).  No continuity or
density argument is involved.

By Lemma~\ref{cp:lem-determinant-quadratic}, the Regge involution fixes
\(x,w,q,s,\mathcal B,\mathcal C\).  It therefore fixes \(r\) and \(r_+\)
and preserves both rational charts.  The two oriented Gram matrices
\((G,w)\) and \((G^R,w)\) are defined over the same chart.

\subsection{The universal open set}
\label{cp:subsec-universal-open}

The rational chart contains points which do not represent nondegenerate
tetrahedra, and it also contains matrices with degenerate coordinate
restrictions.  These unwanted points lie on algebraic divisors and may be
removed once and for all.

Fix one of the two charts in
Proposition~\ref{cp:lem-rational-charts}.

For a regular function \(f\), the \emph{principal open}
\(D(f)\) is obtained by removing its zero set, equivalently by adjoining
\(f^{-1}\) to the coordinate ring.  The \emph{function field} of an
irreducible affine variety is the field of fractions of its coordinate
ring.  Shrinking to a nonempty principal open changes the ring of regular
functions but not this field of rational functions.

\begin{definition}[Universal open]
\label{cp:def-U}
Let \(U\) be the Regge-stable principal open obtained by requiring:
\begin{enumerate}
\item the displayed chart denominator, \(s\), \(w\), and \(1-x^2\) are
      nonzero;
\item \(G\) and \(G^R\) are nonsingular; and
\item every proper principal restriction of \(G\) and \(G^R\) which occurs
      as a coordinate face is nonsingular.
\end{enumerate}
Equivalently, take the product of these finitely many functions and their
Regge transforms and invert that product.  This makes the definition
manifestly Regge stable.  Subsequent constructions may shrink \(U\) further
by finitely many Regge-stable nonvanishing conditions; such a shrinking does
not change its function field.
\end{definition}

Since a rational chart is an open subset of
\(\mathbf G_m^5\times\mathbf A^1\), every nonempty such \(U\) is smooth and
irreducible, and
\begin{equation}
 \mathbf Q(U)=\mathbf Q(a,b,c,\tau,q,r)
 \quad\text{or}\quad
 \mathbf Q(a,b,c,\tau,q,r_+).
 \label{cp:eq-universal-open-function-field}
\end{equation}
The word ``universal'' does not mean that every complex point of \(U\) is a
convex spherical tetrahedron.  It means that \(G\), \(G^R\), their common
orientation root, and the Regge involution are all defined simultaneously
over one algebraic base.  The physical spherical pairs form a distinguished
real locus inside its complex points.

\begin{proposition}[Every physical pair lies on a rational chart]
\label{cp:prop-physical-point-on-chart}
Let \(T,T^R\subset S^3\) be a nondegenerate spherical Regge pair.  Choose
the physical phases \eqref{cp:eq-angle-phases} and the positive root
\(w=\sqrt{\det G(T)}\).  Then the resulting oriented point belongs to at
least one of the two rational charts.  After shrinking that chart as above,
it is a complex point of a universal open \(U\), and its Regge mate is the
image of that point under \(\mathcal R\).
\end{proposition}

\begin{proof}
Because \(0<\theta_{23}<\pi\),
\[
 s=\frac{e^{i\theta_{23}}-e^{-i\theta_{23}}}{2}
   =i\sin\theta_{23}\ne0.
\]
Nondegeneracy and positive definiteness give \(w>0\).  Thus \(sw\ne0\),
and Proposition~\ref{cp:lem-rational-charts} puts the point on
at least one chart.  All principal restrictions of a positive-definite
matrix are positive definite; the same applies to the nondegenerate
spherical mate.  Hence none of the functions removed in the definition of
\(U\) vanishes at the pair.  The common-root lift
\eqref{cp:eq-oriented-regge-lift} identifies the second oriented matrix with
the Regge image of the first.
\end{proof}

This is the step which removes any algebraicity restriction on an individual
tetrahedron.  Its angles may be transcendental: they merely specify a
complex point of the algebraic family \(U\).  The later argument will prove
one identity over \(U\) and then pull that identity back to this point.

\subsection{Recovering the omitted fixed phase}
\label{cp:subsec-fixed-p-cover}

The functions \(a,b,c,q,d=\tau^2/(abc)\) are already defined on \(U\).
The only original angle phase which is not a function on \(U\) is \(p\),
because the rational chart remembers only
\[
 x=-\frac{p+p^{-1}}2.
\]
Equivalently, \(p\) satisfies
\begin{equation}
 p^2+2xp+1=0.
 \label{cp:eq-fixed-p-quadratic}
\end{equation}

\begin{proposition}[The fixed-phase cover]
\label{cp:prop-fixed-p-cover}
Define
\begin{equation}
 \pi:U_p=
 \operatorname{Spec}
 \frac{\mathcal O_U[p]}{(p^2+2xp+1)}
 \longrightarrow U.
 \label{cp:eq-p-cover}
\end{equation}
Then \(\pi\) is finite \'etale of degree two.  Its nontrivial deck
transformation is
\begin{equation}
 p\longmapsto p^{-1}=-2x-p.
 \label{cp:eq-fixed-p-deck}
\end{equation}
The cover is common to \(G\) and \(G^R\), because the Regge transformation
fixes the \(01\)-angle and hence fixes \(x\).
\end{proposition}

\begin{proof}
The polynomial in \eqref{cp:eq-fixed-p-quadratic} is monic of degree two, so
\(U_p\) is finite locally free of rank two over \(U\).  Its discriminant is
\begin{equation}
 (2x)^2-4=4(x^2-1).
 \label{cp:eq-fixed-p-discriminant}
\end{equation}
The function \(1-x^2\) was inverted in the definition of \(U\); hence the
discriminant is a unit and the cover is \'etale.  The product of
the two roots is one, so the second root is \(p^{-1}\), and their sum is
\(-2x\), giving \eqref{cp:eq-fixed-p-deck}.  Finally, \(x\) is fixed by the
Regge involution.
\end{proof}

On a physical fibre the two points of \(U_p\) are
\(e^{i\theta_{01}}\) and \(e^{-i\theta_{01}}\).  The cover may split after
base change, but it need not split over \(U\) itself.  This two-sheeted
fixed-phase datum is the sole residual effect of descending from the full
phase torus to the rational orientation chart.
